\section{Evaluation Design}\label{sec:evaluation}
\subsection{Research questions and evidence levels}\label{sec:rqs}
We evaluate \tool{} with four research questions.
\begin{itemize}
\item \textbf{RQ1 (admission):} Does evidence-bound admission block invalid episodes before production without blocking valid ones?
\item \textbf{RQ2 (cost):} What does validation cost, and is that cost bounded once failed and recovery attempts are charged?
\item \textbf{RQ3 (preservation):} Do permitted operational changes leave the intended scientific experiment unchanged?
\item \textbf{RQ4 (components):} What does each conjunct of Equation~\eqref{eq:ready} contribute beyond a successful smoke run?
\end{itemize}

The evidence comes from three levels of increasing realism and decreasing completeness. Level A consists of executed controller studies on synthetic records. These are a mutation study, a leave-one-conjunct-out ablation, a threshold sensitivity analysis, a bounded-refinement simulation and an assessment-overhead measurement. Level B is a real-GPU integration pilot with small synthetic tensors on the Nibi cluster. Level C is a prospective paired campaign over twelve research artifacts; its availability checks are complete, while installation, adapter validation and comparative reruns are not. Table~\ref{tab:rq-evidence} maps each question to its endpoints and to the evidence currently available. No Level A or B observation is presented as a result of the twelve-artifact campaign.

\begin{table}[t]
\caption{Research questions, endpoints and evidence. ``Executed'' analyses produced the results in Section~\ref{sec:results}; ``prospective'' analyses belong to the campaign protocol in Section~\ref{sec:campaign} and have not been run.}\label{tab:rq-evidence}
\small\begin{tabularx}{\textwidth}{lYYY}
\toprule Question & Endpoint & Executed evidence (Levels A, B) & Prospective evidence (Level C) \\
\midrule
RQ1 & Detection, false blocks, production escapes & Mutation study; threshold sensitivity; four CUDA probe modes & Natural and injected faults across twelve artifacts \\
RQ2 & Allocated cost, time to valid completion & Slurm accounting of the pilot; refinement cost bounds; assessment latency & Paired failure-penalized cost against smoke-then-production \\
RQ3 & Scientific identity, output equivalence & Invariance of $S$ across all refinement proposals & Equivalence tests within predeclared margins \\
RQ4 & Contribution of each conjunct & Leave-one-conjunct-out ablation on the mutation suite & Paired artifact-level ablations \\
\bottomrule\end{tabularx}
\end{table}

\subsection{Implementation details}\label{sec:implementation}
This subsection collects the implementation decisions that realize the approach of Section~\ref{sec:approach}. Table~\ref{tab:implementation} lists the modules, their origin and their size.

\begin{table}[t]
\caption{Prototype and study modules. LOC counts nonblank source lines in the archived revision. ``Inherited'' modules come unchanged from \upstream{} at revision \code{f9dc7d6}; ``Study'' modules generate the reported evidence.}\label{tab:implementation}
\small\begin{tabularx}{\textwidth}{llrY}
\toprule Module & Origin & LOC & Responsibility \\
\midrule
\code{core.py} & New & 130 & Canonical identity, plan validation, readiness (Eq.~\ref{eq:ready}), refinement, discovery \\
\code{pipeline.py} & New & 99 & Input binding, graph/resource correspondence, dry-run or explicit dispatch \\
\code{mcp.py} & New & 52 & Stdio JSON-RPC server exposing three read/propose tools \\
\code{cli.py} & New & 21 & \code{discover}, \code{gate} and \code{refine} commands \\
\code{hpclib.py} & Inherited & 180 & Matrix expansion, dependency resolution and static graph validation \\
\code{run\_stage.py} & Inherited & 184 & Step execution, completion markers and manifests \\
\code{audit.py} & Inherited & 287 & Audit ledger and job watcher \\
\code{collect\_results.py} & Inherited & 282 & Result collection across array cells \\
\code{run\_local.py} & Study & 100 & Level A mutation generator and comparator policies \\
\code{run\_analyses.py} & Study & 176 & Ablation, sensitivity, refinement and overhead analyses \\
\code{probe\_gpu.py} & Study & 70 & Level B CUDA probes with five obligations \\
\bottomrule
\end{tabularx}

\end{table}

\paragraph{Prototype and execution backend.} \tool{} version 0.1.0 is a Python package that requires Python 3.11 or newer. The policy, the command-line interface and the MCP server use only the standard library, which keeps the trusted assessment path free of third-party dependencies. The execution backend reuses the \upstream{} templates pinned at revision \code{f9dc7d6}. Matrix expansion, dependencies, step execution, completion markers, manifests, audit ledgers and result collection are inherited. Source snapshots and full revisions accompany the artifact.

\paragraph{Identity and readiness.} \code{resai.core} computes $h$ as SHA-256 over canonical JSON (sorted keys, compact separators, NaN and infinity rejected), so semantically equal plans hash identically regardless of key order. \code{validate} rejects malformed plans before any evidence is read. These include missing fields, an empty or duplicated probe list, nonpositive budgets, an attempt cap below two, and policy values outside their domains ($0\leq m_{\max}\leq 1$, $0\leq\epsilon\leq 1$, $T\geq 1$\,s). \code{readiness} then evaluates every conjunct of Equation~\eqref{eq:ready} and returns a sorted set of reason codes rather than stopping at the first failure. An operator therefore sees all violated terms at once. Records with nonnumeric or nonfinite telemetry are rejected immediately with an \code{invalid\_telemetry} reason. Table~\ref{tab:operators} lists the reason codes that the conjuncts emit. The prototype trusts operator-supplied records: hashes, timestamps and costs are not tamper-resistant attestations, and the accounting is not a trusted service.

\paragraph{Discovery and refinement.} Discovery (Algorithm~\ref{alg:discover}) issues \code{sinfo --json}, \code{scontrol show partition} and \code{scontrol show topology}, each with a 20-second timeout. It archives return codes and raw output, records an unavailable command as unavailable rather than as an empty cluster, and fixes bandwidth and latency to ``unknown''. It returns raw scheduler observations rather than a normalized node--edge database and implements no topology-aware placement. Metadata reads have per-read timeouts but no global polling limiter; the inherited audit watcher waits at least sixty seconds between updates. \code{refine} implements the single permitted action. It halves a batch size greater than one only when the most recent record matches the current plan, is classified as out of memory, the scientific contract declares batch invariance, and both attempt and compute budgets remain. It deep-copies the plan, modifies only the execution subdocument and returns the proposal without executing it.

\paragraph{Dispatch.} Algorithm~\ref{alg:dispatch} is realized in \code{resai.pipeline}. Declared hashes must bind the selected pipeline's \code{experiments.json} and \code{cluster.env} as well as all other declared inputs; the wrapper checks declared local files but does not discover transitive dependencies. It verifies the resolved graph identity, the scientific metadata declaration, the batch variable and the allocation shape, and rejects raw scheduler overrides in GPU stages. The current mapping supports homogeneous GPU stages with explicit batch configuration and HH:MM:SS walltimes; unsupported shapes block dispatch. Site values containing carriage returns or newlines are rejected. The default is an upstream dry run, and \code{--submit} explicitly invokes submission. Scientific metadata remains an operator assertion about command semantics. Direct upstream invocation bypasses the wrapper, and concurrent changes after verification and duplicate submission remain limitations.

\paragraph{Agent interface.} The stdio MCP server exposes \code{cluster\_snapshot}, \code{assess\_readiness} and \code{propose\_refinement}. It supports initialization, tool listing and tool calls using newline-delimited JSON-RPC for the 2024-11-05 protocol~\cite{mcp2024transport}. File arguments resolve inside the workspace after symlink resolution, and external paths are rejected. No shell or submission tool is exposed. The command-line interface offers the same three operations (\code{discover}, \code{gate}, \code{refine}); \code{gate} exits with status 2 when evidence is not ready, so shell pipelines can stop on a blocked decision. This is a single-user prototype without interoperability certification or a security audit.

\paragraph{Research-artifact adapter.} The MapCoder adapter pins upstream revision \code{c87e6d91} and archives its patch and provenance. A HumanEval-only import patch removes inactive competitive-programming dependencies. A local Transformers backend replaces the cloud backend, while the selected strategy's prompts and traversal are retained. Generated programs run in an AST-restricted subprocess with CPU, address-space and wall-time limits, rather than the original thread evaluator, and assertions remain enabled. This narrow evaluator is not a hardened sandbox. Per-attempt directories prevent old output files from appearing fresh. The resume probe reloads persisted results into fresh strategy objects and checks that no additional inference is needed. Requirements pin PyTorch 2.14.1, Transformers 4.57.6 and tiktoken 0.12.0.

\paragraph{Analysis pipeline and tests.} CSV and JSON records retain attempt identifiers, source hashes, software and device observations, and Slurm accounting. Each generator asserts the expected record cardinality before producing a table or plot; for example, the ablation's full variant must reproduce the mutation study's \tool{} column exactly. All figures are vector PDFs generated from the archived records. Thirteen unit tests cover policy edge cases, path confinement and graph/resource mismatches, and a claims checker verifies that reported counts, citations and frozen identities match the raw data.

\subsection{Experimental setup}\label{sec:experimental-setup}
\paragraph{Environments.} Level A analyses ran on a Windows 11 workstation (Intel x86-64) with Python \AnalysisPython{}. Level B records report PyTorch 2.14.1 and CUDA runtime 13.2 on an H100 1g.10gb MIG instance of Nibi. CUDA execution and environment preparation take place inside scheduled allocations. MIG-instance seconds and full-H100 seconds are different resource units and are never pooled. Table~\ref{tab:setup} summarizes the workloads.

\begin{table}[t]
\caption{Workloads and execution settings. MapCoder entries describe the frozen requested configuration; its execution status is reported separately. CPU and RAM values are per allocation.}\label{tab:setup}
\small\begin{tabularx}{\textwidth}{lYYY}
\toprule Factor & Controller studies (A) & CUDA integration (B) & Staged MapCoder \\
\midrule
Workload & \FaultFamilies{} fault families + clean control; parameter sweeps; refinement episodes & MLP 512--1024--64; input $128\times512$ & HumanEval/0; Qwen2.5-Coder-7B-Instruct \\
Repetitions & Seeds 0--19 per cell; \OverheadRepeats{} timing repeats & Two cells; four modes each & Two fresh sequential model processes \\
Accelerator & None & One H100 1g.10gb MIG per cell & One full H100 requested \\
Host resources & Local workstation & 2 CPU cores; 12 GB RAM & 4 CPU cores; 48 GB RAM \\
Allocation ceiling & No Slurm allocation & 15 min per cell; 30 instance-min total & 15 min for both processes \\
Output oracle & Injected property / clean control & Probe coverage and output validity & Record validity; task correctness separate \\
\bottomrule\end{tabularx}
\end{table}

\paragraph{Policy parameters.} Table~\ref{tab:policy-settings} lists the readiness parameters, fixed before analysis. They exercise the contract rather than representing optimized thresholds; the sensitivity analysis below varies them explicitly. A production study must justify its own envelope and should distinguish framework-allocated memory from all memory used on a device.

\begin{table}[t]
\caption{Readiness parameters fixed before analysis. Process-instance budgets govern the supplied evidence; Slurm ceilings separately bound allocations.}\label{tab:policy-settings}
\small\begin{tabularx}{\textwidth}{Yrrr}
\toprule Parameter & Symbol & Controller / CUDA & MapCoder \\
\midrule
Matching observations required & --- & 2 & 2 \\
Evidence lifetime (s) & $T$ & 3,600 & 3,600 \\
Maximum memory fraction & $m_{\max}$ & 0.90 & 0.90 \\
Maximum between-attempt difference & $\epsilon$ & 0.05 & 0.05 \\
Maximum supplied attempts & $K$ & 4 & 4 \\
Process-instance validation budget (s) & $C$ & 1,200 & 1,800 \\
\bottomrule\end{tabularx}
\end{table}

\paragraph{Mutation study (RQ1).} A generator creates a valid plan with two successful records, varying a data-identity token and a baseline memory fraction drawn uniformly from $[0.4,0.7]$ under \LocalSeeds{} deterministic seeds. For each seed, one of \FaultFamilies{} operators corrupts a single property, and a clean control is retained. Table~\ref{tab:operators} defines every operator, the term of Equation~\eqref{eq:ready} it violates and the reason codes it produces. This yields \CasesPerCondition{} cases per policy and \LocalDecisions{} decisions overall. Four policies are compared. \emph{Direct} admission accepts everything and serves as a reference floor. \emph{Static} applies the actual upstream graph validator. The \emph{smoke proxy} combines that validator with successful process statuses and output-file existence. It models what a passing smoke run shows, but it does not execute the entire upstream smoke pipeline or every check an experienced user would perform. Finally, \tool{} combines upstream validation with Equation~\eqref{eq:ready}. The oracle follows the operator: a clean fixture should pass and an intentionally violated property should block.

\begin{table}[t]
\caption{Fault operators of the mutation study. Each operator corrupts one property of an otherwise valid plan and evidence pair. Component is the ablation group that owns the violated term (Graph faults are caught by the upstream validator). First blocker is the weakest policy, in the order static, smoke proxy, \tool{}, that blocks all \LocalSeeds{} variants; each stronger policy blocks them too, and direct admission blocks none.}\label{tab:operators}
\footnotesize\setlength{\tabcolsep}{3pt}\begin{tabularx}{\textwidth}{>{\raggedright\arraybackslash}p{2.1cm}>{\raggedright\arraybackslash}Yll>{\raggedright\arraybackslash}p{2.9cm}l}
\toprule Fault family & Mutation & Violated term & Component & Emitted reason codes & First blocker \\
\midrule
GPU request absent & Remove the stage GPU request & $\operatorname{Valid}(P)$ & Graph & --- (graph validator) & Static \\
Dependency invalid & Depend on a nonexistent stage & $\operatorname{Valid}(P)$ & Graph & --- (graph validator) & Static \\
Process failed & Set last status to failed & $\operatorname{OK}(e_i)$ & Outcome & \code{execution\_\allowbreak{}failure} & Smoke proxy \\
Required probe absent & Drop the checkpoint obligation & $O\subseteq C_i$ & Outcome & \code{coverage} & \tool{} \\
Output invalid & Mark the present output invalid & $\operatorname{OutputValid}(e_i)$ & Outcome & \code{output\_\allowbreak{}contract} & \tool{} \\
Plan identity changed & Change the scientific model & $h_i=h$ & Identity & \code{two\_\allowbreak{}distinct\_\allowbreak{}attempts} & \tool{} \\
Attempt ID duplicated & Reuse the first attempt ID & distinct IDs & Identity & \code{duplicate\_\allowbreak{}or\_\allowbreak{}missing\_\allowbreak{}attempt\_\allowbreak{}id}, \code{two\_\allowbreak{}distinct\_\allowbreak{}attempts} & \tool{} \\
Evidence expired & Move the last record to $t=1$ & $t-t_i\leq T$ & Freshness & \code{stale\_\allowbreak{}evidence} & \tool{} \\
Memory limit exceeded & Set peak memory to 0.98 (OOM) & $m_i\leq m_{\max}$ & Memory & \code{memory\_\allowbreak{}headroom}, \code{unstable\_\allowbreak{}memory} & \tool{} \\
Memory unstable & Add 0.12 to the last peak & $\Delta_{\mathrm{smoke}}\leq\epsilon$ & Memory & \code{unstable\_\allowbreak{}memory} & \tool{} \\
Validation budget exceeded & Charge 1,300 s to the last record & $\sum c_i\leq C$ & Budget & \code{compute\_\allowbreak{}budget} & \tool{} \\
Accounting absent & Delete the last cost $c_i$ & finite $c_i>0$ & Budget & \code{invalid\_\allowbreak{}telemetry:gpu\_\allowbreak{}seconds} & \tool{} \\
\bottomrule
\end{tabularx}

\end{table}

\paragraph{Ablation (RQ4).} Because Equation~\eqref{eq:ready} is a conjunction, removing a conjunct is equivalent to ignoring the reason codes it emits. We group reason codes into five components: \emph{Identity} ($h_i=h$, distinct identifiers), \emph{Outcome} ($\operatorname{OK}$, coverage, output validity), \emph{Freshness} ($t-t_i\leq T$), \emph{Memory} (headroom and stability) and \emph{Budget} (attempt cap, compute cap and valid accounting). A variant admits a case when the upstream validator passes and every emitted reason belongs to its removed component. All \AblationVariants{} variants are evaluated on the mutation suite, holding the fixtures, seeds and validator fixed. The ablation measures the logical contribution of each conjunct on constructed faults, not the frequency of the corresponding faults in practice.

\paragraph{Threshold sensitivity (RQ1).} To test whether decisions follow the declared parameters, we sweep one evidence property at a time across three settings of the parameter that governs it, holding all else valid. The swept properties are the age of both matching records ($T\in\{1800,3600,7200\}$\,s, 250\,s steps), a footprint scale applied to each seed's baseline memory ($m_{\max}\in\{0.8,0.9,0.95\}$), the between-attempt memory difference ($\epsilon\in\{0.02,0.05,0.1\}$) and the total charged validation time ($C\in\{600,1200,1800\}$\,s). Each point is assessed for \LocalSeeds{} seeds, giving \SensitivityAssessments{} assessments over \SensitivitySettings{} parameter settings. Plans are rehashed after every change so that only the swept property differs.

\paragraph{Bounded refinement (RQ2, RQ3).} We simulate Algorithm~\ref{alg:smoke} with an operator who adopts every proposal. Peak memory scales linearly with batch size from each seed's baseline at batch 8. A probe whose peak exceeds device capacity fails with an out-of-memory classification, and every probe is charged 60\,s. Episodes start from batch sizes $\{8,16,32,64,128\}$ under attempt caps $K\in\{2,\dots,8\}$, with and without declared batch invariance, for \LocalSeeds{} seeds each: \RefinementEpisodes{} episodes. An episode ends READY, BLOCKED with no applicable proposal, or BLOCKED because the attempt or compute budget is exhausted. The scientific subdocument is hashed before and after every proposal. The memory model is a deliberate simplification that exercises the control loop; it does not predict the memory behavior of real workloads.

\paragraph{Assessment overhead (RQ2).} We time a single \code{readiness} call for evidence ledgers of $2^1$ to $2^{11}$ matching records, \OverheadRepeats{} repetitions each, using a monotonic high-resolution clock. We report the median and interquartile range and compare them with the shortest measured Level B probe allocation.

\paragraph{Integration pilot (RQ1, RQ2).} The Level B pilot submits a two-cell Slurm array through \upstream{}, followed by CPU analysis and collection under a CPU association. Each cell runs four modes. Clean mode checks imports, initialization, input execution, output cardinality and checkpoint round trip. The other modes remove an output row, omit checkpoint testing, or raise a controlled runtime exception after CUDA execution. Failures are recorded without terminating the enclosing collection stage, so a process-level failure cannot silently remove an inconvenient observation. Analysis requires exactly four modes by two repeats, pairs repeats within each mode and applies Equation~\eqref{eq:ready}. A per-mode plan identity is assigned retrospectively to describe the known probe; this is not pre-execution admission of the full graph. Slurm allocation and queue times are reported separately from timed probe computation, which excludes installation, imports and scheduler overhead. This is synthetic integration, not an agent-paper rerun or a large-model benchmark.

\paragraph{Staged MapCoder pilot (RQ3).} Before expanding to twelve artifacts, we apply the approach to the local-model MapCoder configuration~\cite{resaiartifact7}. The specification fixes a cached Qwen revision (\code{c03e6d35}), greedy decoding with seed zero, three recalled examples, five debugging opportunities and one outer attempt. Input and generation caps are 12,288 and 4,096 tokens, batch size is one, and batch invariance is false. Task, model revision, prompts and scoring belong to $S$. Consequently, the local-model adapter is a new deployment configuration, not a reproduction of published aggregate accuracy. The protocol calls for two fresh model processes executed sequentially in a scheduled interactive smoke allocation. Obligations cover imports, model initialization, task execution, output schema and cardinality, and persisted-result resume. Scientific task correctness is recorded separately: a valid record of an unsuccessful generated program is not a malformed output. We require an initial blocked assessment without evidence, then assess collected observations; a dependent CPU job collects evidence and prints a dispatch dry run. No larger campaign is dispatched before author review.

\paragraph{Analysis of executed levels.} Level A cases are deterministic functions of their seeds, and seed variants are not independent research artifacts. We therefore report counts, fractions and per-family matrices without significance tests or population confidence intervals. Operators were authored with knowledge of the policy, so a perfect count establishes conformance to the specification, not accuracy on naturally occurring faults. Level B reports all eight raw records and the corresponding Slurm accounting.

\subsection{Prospective twelve-artifact campaign}\label{sec:campaign}
\paragraph{Hypotheses.} The campaign tests four primary hypotheses at the artifact level; all are two-sided except the equivalence hypothesis. $H_{01}$ states zero paired difference in pre-production detection rate between full \tool{} and the strongest smoke comparator, with false-block rates reported alongside. $H_{02}$ states zero paired difference in the predeclared failure-penalized time/resource endpoint. For a declared output difference $d$ and equivalence margin $\delta$, $H_{03}:d\leq-\delta\;\lor\;d\geq\delta$ is tested using two one-sided tests against $-\delta<d<\delta$; equivalence must be established, not inferred from a nonsignificant difference. $H_{04}$ states zero paired difference between the full policy and each predeclared ablation on its targeted outcome. None of these hypotheses is tested with the synthetic seed variants of Level A.

\paragraph{Subjects.} We assembled twelve real agent research artifacts spanning repository repair, generation, collaborative development and test-environment construction (Table~\ref{tab:subjects}). Paper links and full revisions are in \code{experiments/subjects.json}. Availability was checked through the public repository API on 23 October 2026; a current revision is an availability observation, not automatically the paper's evaluated version.
\begin{table}[t]
\caption{Candidate artifacts of the prospective campaign. Every comparative execution status is currently \emph{not run}. HEADs are availability observations, not validated paper-version pins.}\label{tab:subjects}
\small\begin{tabularx}{\textwidth}{lXl}
\toprule Artifact & Intended workflow & Observed HEAD \\
\midrule
\href{https://github.com/SWE-agent/SWE-agent}{SWE-agent}~\cite{yang2024sweagent} & repository repair & \texttt{3ea751c0} \\
\href{https://github.com/OpenHands/OpenHands}{OpenHands}~\cite{wang2025openhands} & software development & \texttt{11eba891} \\
\href{https://github.com/AutoCodeRoverSG/auto-code-rover}{AutoCodeRover}~\cite{resaiartifact3} & structure-aware repair & \texttt{585d3e63} \\
\href{https://github.com/FoundationAgents/MetaGPT}{MetaGPT}~\cite{resaiartifact4} & role-based development & \texttt{11cdf466} \\
\href{https://github.com/OpenBMB/ChatDev}{ChatDev}~\cite{resaiartifact5} & communicative development & \texttt{a2ea4a26} \\
\href{https://github.com/huangd1999/AgentCoder}{AgentCoder}~\cite{resaiartifact6} & generation and testing & \texttt{d8538d63} \\
\href{https://github.com/Md-Ashraful-Pramanik/MapCoder}{MapCoder}~\cite{resaiartifact7} & competitive programming & \texttt{c87e6d91} \\
\href{https://github.com/sola-st/RepairAgent}{RepairAgent}~\cite{resaiartifact8} & Java repair & \texttt{5d9dc8de} \\
\href{https://github.com/sola-st/ExecutionAgent}{ExecutionAgent}~\cite{bouzenia2025executionagent} & test environment setup & \texttt{8fb16de8} \\
\href{https://github.com/xingyaoww/code-act}{CodeAct}~\cite{resaiartifact10} & executable code actions & \texttt{d607f56c} \\
\href{https://github.com/FSoft-AI4Code/HyperAgent}{HyperAgent}~\cite{resaiartifact11} & generalist SE & \texttt{c3092f26} \\
\href{https://github.com/FSoft-AI4Code/AgileCoder}{AgileCoder}~\cite{resaiartifact12} & sprint-based development & \texttt{f9924185} \\
\bottomrule
\end{tabularx}

\end{table}

The inventory is purposive, not a systematic literature-review funnel. Before claiming generality, the campaign must archive searches of ICSE, FSE, ASE, ISSTA, EMSE, TSE and relevant AI venues for 2023--2026, combining ``agent'', ``software engineering'', ``code generation'', ``program repair'' and ``test execution''. Screening must retain identifiers, duplicate decisions, queries, paper and repository links, and exclusions. Eligibility requires accessible paper-compatible code, permission for the experiment, a reduced executable workload, an output oracle and deployment compatible with site policy. Infeasible candidates remain in the ledger, and replacements must be recorded before comparative outcomes are inspected.

The set is concentrated in generation and repair. CodeAct includes tasks beyond software engineering, so its inclusion requires restricting evaluation to software-oriented tasks. Cloud-API and local-model workloads must be separated: API experiments can establish deployment and CPU/allocation observations but cannot establish local GPU savings. A local-endpoint port is a new configuration requiring capability checks and pinned weights. Container-dependent artifacts require a site-approved alternative rather than privileged Docker on compute nodes.

\paragraph{Procedure.} For each eligible artifact, we freeze paper-compatible source, environment, benchmark subset, model, serving configuration, prompts, token limits, seeds and scoring, and specify the resource shape and permitted repairs before running any condition. A held-out validation subset is used to design probes; final benchmark outcomes are not fed into the admission policy. Native retries, recovery and checkpoints remain enabled in every condition to avoid handicapping comparators. The conditions are direct native submission, upstream static validation, upstream smoke-then-production, full \tool{}, and three ablations removing coverage obligations, identity/freshness checks and refinement proposals. All conditions share the artifact adapter. The design crosses twelve artifacts, seven conditions, five paired seeds and four episode classes (clean, environment fault, resource fault and output-contract fault): 1,680 capped episodes if all artifacts are eligible. This is a design count, not submitted work or a cost estimate, and a feasibility pilot must fix per-artifact caps and the total budget before launch. Conditions are randomized within artifact, site/resource and submission-time blocks. Caches are held constant or varied as a declared factor, and documented faults are injected at the same point in paired conditions. Natural faults are retained separately because their frequency and detectability may differ from seeded faults. Human intervention, transcript identity, repair actions, completion, validated output and allocation identifiers are logged. Failed installations and time-capped episodes remain in the denominator.

\paragraph{Metrics.} RQ1 measures detection among faulty episodes and false blocking among valid clean episodes, alongside production escapes and blocking reasons. Rejecting everything yields high detection and unusable availability, so both outcomes are mandatory. RQ2 charges all validation, recovery and production allocations:
\begin{equation}
C_{\mathrm{total}}=\sum_{j\in\mathrm{validation,recovery,production}}a_j\operatorname{ElapsedAllocated}(j),
\end{equation}
where $a_j$ counts allocated accelerator instances. CPU-seconds, instance type, API token cost and queue time are reported separately, and unlike accelerators are never pooled into unlabeled GPU-hours. Early rejection may avoid wasted work while preventing a valid result, so completion rate and time-to-result are needed to interpret resource use. Each episode receives a predeclared cap, and actual allocated consumption is always retained. A separate failure-penalized endpoint assigns the cap to unsuccessful episodes for the primary intention-to-run comparison; this imputation is not labeled measured consumption. A sensitivity analysis uses restricted mean time to valid completion under explicit censoring rules. Queue delay is submission-to-start, not allocated GPU time. RQ3 requires exact immutable-input identity and artifact-specific output tolerances fixed before comparison: deterministic outputs may require equality, whereas stochastic task success needs equivalence margins and sufficient repetitions. RQ4 compares the full policy against each ablation on targeted outcomes, including extra validation cost, and effects must persist against the full upstream smoke condition, not merely direct submission.

\paragraph{Statistical analysis.} The primary unit is the artifact; seeds, tasks and injected faults are nested observations. We compute artifact-level paired differences, report their distribution with a bootstrap interval that resamples artifacts, and perform a paired randomization test over artifact-level signs. With twelve eligible artifacts, all $2^{12}$ sign configurations can be enumerated. Holm correction is applied within the predeclared primary family, and effect sizes are reported with their uncertainty. Twelve artifacts provide limited precision, so a prospective precision or power assessment is required before fixing the campaign size~\cite{wohlin2012experimentation}.
